\paragraph{4.1. Задача поиска максимального подмассива}

\subparagraph{4.1.1}
Что возвращает процедура \textsc{Find-Maximum-Subarray}, когда все элементы \(A\) отрицательны?

\textbf{Ответ:}
Так как основную работу выполняет функция \textsc{Find-Max-Crossing-Subarray}, можно заметить что начальными значениями в алгоритме явялется $-\infty$
следовательно алгоритм вернёт максимальную по сумме подмассив, а точнее массив из одного элемента с максимальным значением. 

\subparagraph{4.1.2}
Напишите псевдокод для решения задачи поиска максимального подмассива методом грубой силы.
Ваша процедура должна выполняться за время \(\Theta(n^2)\).

\textbf{Ответ:}

\begin{lstlisting}
std::tuple<size_t, size_t, long> FindMaximumSubarray(const std::vector<int>& A)
{
    size_t begin = 0;
    size_t end = 0;
    long maxSum = std::numeric_limits<long>::min();

    for(size_t i = 0; i < A.size(); ++i)
    {
        long localSum = 0;
        for (size_t y = i; y < A.size(); ++y)
        {
            localSum += A[y];
            if (localSum > maxSum)
            {
                begin = i;
                end = y;
                maxSum = localSum;
            }
        }
    }

    return {begin, end, maxSum};
}
\end{lstlisting}

\subparagraph{4.1.3}
Реализуйте и метод грубой силы, и рекурсивный алгоритм на своем компьютере.
Каким оказывается размер задачи точки пересечения \(n_0\),
в которой рекурсивный алгоритм превосходит алгоритм грубой силы?
Далее измените базовый случай рекуррентного алгоритма,
применяя алгоритм грубой силы при размере задачи, не превосходящем \(n_0\).
Меняет ли это точку пересечения?

\textbf{Ответ:}

Сделаю чуть позже. В бенчмарках.

\subparagraph{4.1.4}
Предположим, что мы меняем определение задачи поиска максимального подмассива,
позволяя конечному результату быть пустым массивом и полагая,
что сумма значений пустого массива равна нулю. 
Как бы вы изменили любой алгоритм, не допускающий решения в виде пустого массива, 
чтобы такой результат в виде пустого подмассива стал возможным?

\textbf{Ответ:}

\begin{itemize}
    \item Инициализировать $\text{max-sum} = 0$ вместо $-\infty$.
    \item Инициализировать $\text{left-sum} = 0$ и $\text{right-sum} = 0$ в функции 
        \textsc{Find-Max-Crossing-Subarray}.
    \item Если итоговая максимальная сумма равна $0$, вернуть пустой подмассив 
        (например, с индексами $(-1, -1)$).
\end{itemize}